\documentclass[11pt]{article}
\usepackage[margin=1in]{geometry}
\usepackage{amsmath,amssymb,amsthm,mathtools}
\usepackage[hidelinks]{hyperref}
\newtheorem{theorem}{Theorem}
\newtheorem{lemma}[theorem]{Lemma}
\newtheorem{corollary}[theorem]{Corollary}
\newcommand{\Q}{\mathbb Q}
\newcommand{\C}{\mathbb C}
\newcommand{\F}{\mathbb F}
\newcommand{\Z}{\mathbb Z}
\title{A uniform large-exponent coefficient criterion}
\author{User-supplied coefficient argument checked and integrated; retained prime-power proof credited}
\date{17 November 2026}
\begin{document}
\maketitle
This note supplies the elementary large-exponent criterion and prime-power input used by the completed \href{ALL_UNIVARIATE_TWO_TERM_SFC.tex}{all-index theorem}. It makes no separate exceptional-set claim.

\section{A coefficient estimate excluding every repeated root}
\begin{lemma}\label{lem:factorial}
For integers $1\le u\le v$,
\[
 \frac{u!}{(v!)^{u/v}}\le\left(\frac uv\right)^{u/2}. \tag{1}
\]
\end{lemma}
\begin{proof}
The sequence $a_k=(k!)^{1/k}/\sqrt k$ is nondecreasing.
Indeed AM--GM gives $k!\le((k+1)/2)^k$, while
$(1+1/k)^{k+1}\le4$: the function
$(x+1)\log(1+1/x)$ decreases for $x>0$, since its derivative is
$\log(1+1/x)-1/x<0$, and at $x=1$ its exponential is four.
Together these imply
\[
 (k!)^2(k+1)^{k(k-1)}\le k^{k(k+1)},
\]
which, after taking the appropriate positive powers, is exactly
$a_k\le a_{k+1}$. Comparing $a_u$ and $a_v$ proves (1).
\end{proof}

For $1\le j<n$ we also have
\[
 j\log(n/j)\ge(\log2)\min(j,n-j). \tag{2}
\]
For $j\le n/2$ this follows from $n/j\ge2$.
Otherwise use concavity of $-x\log x$ above its chord between
$x=1/2$ and $x=1$, then multiply by $n$.

\begin{lemma}[Two weighted coefficient masses]\label{lem:masses}
For $P(z)=1+z^n+\sum_{j=1}^{n-1}b_jz^j$, put
\[
 A=\sum_{j=1}^{n-1}\frac jn|b_j|,\qquad
 B=\sum_{j=1}^{n-1}\left(1-\frac jn\right)|b_j|.
\]
If $A<1$ and $B<1$, $P$ has no repeated complex root.
\end{lemma}
\begin{proof}
If $P(\alpha)=P'(\alpha)=0$ and $|\alpha|\le1$, subtracting
$(\alpha/n)P'(\alpha)$ from $P(\alpha)$ yields $1\le B$, a contradiction.
If $|\alpha|\ge1$, divide $(\alpha/n)P'(\alpha)=0$ by $\alpha^n$.
The intermediate powers have modulus at most one, giving $1\le A$.
The constant term already excludes $\alpha=0$.
\end{proof}

\begin{theorem}\label{thm:large}
For every pair of positive integers $n,m$,
\[
 \boxed{2^m\ge4n^2\quad\Longrightarrow\quad F_{n,m}\text{ is squarefree}.}
\]
The conclusion holds over $\C$, and hence over $\Q$.
\end{theorem}
\begin{proof}
Set $r=((mn)!)^{-1/n}>0$ and normalize
\[
 P(z)=F_{n,m}(rz)=1+z^n+\sum_{j=1}^{n-1}b_jz^j.
\]
By (1) at $u=mj,v=mn$ and then (2), with $q=2^{-m/2}$,
\[
 0<b_j\le\binom nj(j/n)^{mj/2}
 \le c_j:=\binom nj q^{\min(j,n-j)},\qquad c_j=c_{n-j}.
\]
Symmetry gives $A,B\le\frac12\sum_{j=1}^{n-1}c_j$.
Since $q^{\min(j,n-j)}\le q^j+q^{n-j}$,
\[
 A,B\le\frac12\sum_{j=1}^{n-1}c_j\le(1+q)^n-1.
\]
Under the displayed criterion, $q\le1/(2n)$, so
\[
 A,B\le\left(1+\frac1{2n}\right)^n-1
 <e^{1/2}-1<1.
\]
The final inequality follows already from $e<4$.
Lemma~\ref{lem:masses} excludes all repeated roots, and the nonzero scaling
preserves their multiplicities. If $n=1$, the sums are empty and the
polynomial is directly linear. Equality in $2^m\ge4n^2$ is allowed.
\end{proof}

\section{The retained prime-power theorem and the exponent one}
The independent pointwise result already follows from
\href{PRIME_POWER_TWO_TERM_SFC.tex}{the registered unit-deformation proof}:
for every $n\ge1$, $F_{n,m}$ is squarefree whenever $m=1$ or $m=p^r$.
Here is its precise match to the proposed formulas.
For $r\ge0$, set $m=p^r$ and $c=(p^r-1)/(p-1)$; $r=0$ means $m=1,c=0$.
After $T=t/p^c$, the coefficient is
\[
 \binom nk\frac{(p^rk)!}{p^{ck}}
 =\frac{n!}{(n-k)!}W_k,\qquad
 W_k=\frac{(p^rk)!}{p^{ck}k!}\in\Z_p^\times,
 \quad W_k\equiv(-1)^{ck}\pmod p.
\]
The unit congruence and digit hull are credited to the pinned
\href{sources/rocks-2015-author-preprint.pdf}{Rocks 2015 author manuscript},
Lemma~3.8 and Proposition~3.22, and rederived in the retained proof.
The new note uses that source as an author manuscript, not an unavailable
final journal version or a priority assertion.

For $n=\sum_s d_sp^s$, a nonzero digit $a=d_s$ gives $K=\sum_{\ell<s}d_\ell p^\ell$,
$e=p^s$ and $h=(e-1)/(p-1)$.
Its normalized edge residual, with constant coefficient one, is exactly
\[
 R_s(Y)=\sum_{j=0}^a\frac{a!}{(a-j)!}(\sigma_sY)^j,
 \qquad\sigma_s=(-1)^{h+ce}. \tag{3}
\]
This is the existing formula $\epsilon_s=(-1)^h\gamma^e$ with
$\gamma=(-1)^c$. Wilson's factorial-unit identity
$U_p(N!)\equiv(-1)^{v_p(N!)}\prod_\ell N_\ell!$ gives the coefficient ratios.
For $a<p$, the truncated exponential $E_a(Z)=\sum_{j=0}^aZ^j/j!$
satisfies $E_a'=E_{a-1}$ and $E_a-E_a'=Z^a/a!$.
A common root would be zero, but $E_a(0)=1$.
Thus (3), the reversed and rescaled $a!Y^aE_a(1/Y)$, is squarefree.
The retained weighted-initial multiplication lemma excludes every
nonconstant square factor over $\Q_p$.
This tests the residual in $Y$, before substituting a power $Y=Z^e$;
it assumes no separability of that substitution in characteristic $p$.
For $m=1$, take any prime, for example $p=2$, and $W_k=1$.
The same deformation theorem therefore makes the exponent-one case explicit.

\begin{corollary}\label{cor:candidates}
Every hypothetical nonsquarefree pair satisfies
\[
 \boxed{m<2\log_2(2n),\qquad m\text{ has at least two distinct prime divisors}.}
\]
For fixed $n$, these are a finite list of necessary candidate exponents,
not assertions that any candidate actually fails squarefreeness.
Other registered index certificates may already eliminate individual
entries in this necessary list; this filter does not reopen those results.
\end{corollary}
\begin{proof}
Negate Theorem~\ref{thm:large} and exclude $m=1$ and all prime powers.
The remaining positive integers are precisely those with at least two
distinct prime divisors.
\end{proof}
This coefficient criterion alone leaves exponent six in the possible
range for $n>4$. The separate
\href{ALL_UNIVARIATE_TWO_TERM_SFC.tex}{all-index theorem} excludes those
remaining pairs throughout the ordinary constant-plus-power family.


\end{document}
